% !TEX root = main.tex

\section{Conclusion and Recommendations}
\label{sec:conclusion}

This study asked where the predictability in daily large-capitalisation
technology equity returns resides, and answered it by decomposing a predictive
distribution into the parts a forecaster can actually control.

The conditional mean is not one of them. Across eight assets, six model classes
and two evaluation protocols, nothing improves on a per-ticker historical
average. Zero of 48 model--ticker tests survive multiplicity correction, and the
count of uncorrected rejections falls below what chance alone produces. When
four sequence architectures are permitted to choose their own regularisation on
held-out data, all four choose to discard the predictors entirely and forecast a
constant. The result is not that flexible models fail to find the signal; it is
that, offered a fair choice, they decline to look for it.

The conditional variance is a different matter. All three GARCH-family
specifications beat a rolling-variance benchmark at $p<0.001$, and the
conditional variance accounts for $54\%$ of the total gain in log predictive
density over a null forecast. A further $45\%$ comes from the shape of the
predictive distribution---Student-$t$ innovations and a scale calibrated on
validation data. The conditional mean accounts for $0.7\%$.

\subsection{Recommendations for practice}

\textbf{Spend modelling effort on the second moment.} At the daily horizon the
return of effort in conditional-mean modelling is close to zero, while
conditional-variance modelling yields large and testable gains. A practitioner
choosing where to allocate attention should not divide it evenly.

\textbf{Report calibration, not only accuracy.} Every conclusion above that
matters is invisible to RMSE. The four rungs of
Section~\ref{sec:decomposition} differ by less than $0.4\%$ in RMSE and by
$0.28$ nats in log score. A study reporting only point accuracy would conclude
that nothing distinguishes them.

\textbf{Test differences rather than ranking them.} Several apparent orderings
in this study dissolve under testing: EGARCH's QLIKE advantage over GARCH
($p=0.33$), the single per-ticker rejection that fails multiplicity correction,
and---in an earlier version of this pipeline---an entire ranking of volatility
models produced by a diverged fit.

\textbf{Separate statistical from economic significance.} The conditional mean's
contribution is real at $p=0.029$ and negligible at $0.7\%$ of the total. With
several thousand observations these two facts coexist routinely, and reporting
only the first would misrepresent the finding.

\textbf{Declare features; never infer them from column type.} The single most
consequential defect in the earlier version of this work was a feature-selection
rule that admitted any numeric column. It produced no error, no warning and no
implausible training diagnostic, and it invalidated every out-of-sample number
in the study.

\subsection{Limitations}

The scope is eight assets in one sector over one decade. The between-ticker
variance is identified by eight groups, which is few, and the transfer result of
Section~\ref{sec:transfer} should be read as evidence about large-capitalisation
technology equities rather than about equities generally.

GARCH parameters are estimated once on the training block and then held fixed.
This is conservative and avoids look-ahead, but Section~\ref{sec:calibration-results}
shows it also understates later volatility by roughly twelve percent. Rolling
re-estimation is the obvious extension and is not attempted here.

Calibration is improved but not achieved: PIT uniformity is still rejected at
$p=0.012$. Candidate explanations include remaining asymmetry in the conditional
distribution, time variation in the degrees of freedom, and the single scalar
used for scale calibration where a time-varying factor may be warranted.

The realised-variance proxy is the squared daily return. QLIKE is robust to
proxy noise \citep{patton2011volatility}, but intraday data would support sharper
volatility comparisons than are possible here.

Finally, the study conditions on one information set. The absence of
conditional-mean predictability is a statement about these fourteen predictors
over this window, and is consistent with, rather than a demonstration of, the
episodic predictability described by \citet{timmermann2008elusive}.

\subsection{Extensions}

Three directions follow directly. Rolling GARCH re-estimation would test whether
the scale miscalibration is an artefact of frozen parameters, and would make the
calibration step unnecessary if so. A skewed or time-varying-shape predictive
distribution would address the residual PIT non-uniformity. And extending the
panel across sectors would test whether the near-exchangeability found here is a
property of large-capitalisation technology equities specifically, or a more
general feature of the asset class.

\subsection{Reproducibility}

All data are public: prices from Yahoo Finance, macro-financial series from the
Federal Reserve Economic Data service. The full pipeline---panel construction,
feature dictionary, volatility estimation, the Bayesian sampler, both evaluation
protocols, all scoring rules and every figure and table in this
paper---is released with pinned dependencies. Two audit scripts are included
that exit non-zero when they detect the predictor-drift and prediction-scale
failures described in Section~\ref{sec:extrapolation}, so that the defect which
invalidated the earlier analysis cannot recur silently.
